% !TEX program = xelatex
% Bulk-vs-tail across all three arms: does the plain-CRPS AR(2) effect survive
% tail restriction? Table generated by make_bulk_tail_table.R.
\documentclass[12pt]{article}

\usepackage{fontspec}
\usepackage{xeCJK}
\setCJKmainfont[AutoFakeSlant=.2, AutoFakeBold=3]{AR PL KaitiM Big5}

\usepackage[margin=1in]{geometry}
\usepackage{amsmath, amssymb}
\usepackage{booktabs}
\usepackage{array}
\usepackage{hyperref}
\hypersetup{colorlinks=true, linkcolor=blue!45!black, urlcolor=blue!45!black}

\title{\textbf{bulk 還是 tail？三個 arm 的統一裁決}}
\author{把 twCRPS 補到模擬與 time-block}
\date{2026 年 7 月}

\begin{document}
\maketitle

\section*{問題}

臭氧上我們發現 plain CRPS 的「AR(2) 較差」是 bulk 效應（尾部的 Brier 與 twCRPS
都測不出）。本文把同樣的\textbf{尾部帶 twCRPS}（$\int_{q_{.90}}^{q_{.99}}
\mathrm{BS}(z)\,dz$，$z$-空間梯形積分）補到\textbf{模擬空間 hold-out}與
\textbf{time-block forecast}兩個 arm，檢查 bulk-vs-tail 是否普遍成立。每個 arm
並列三個分數：兩個尾部焦點分數（單門檻 Brier、尾部帶 twCRPS）對上 bulk 主導的
plain CRPS。

\bigskip
% Auto-generated by make_bulk_tail_table.R -- do not edit by hand.
\begin{table}[htbp]
\centering
\caption{Bulk-vs-tail across the three arms for the AR(2) contrast. Two
tail-focused scores (single-threshold Brier, tail-band twCRPS) vs the
bulk-dominated plain CRPS; verdict by each row's own test at $\alpha=0.05$
(ozone: two-sided site-clustered bootstrap $p$; simulation and time-block:
one-sided Wilcoxon $p_W$ over $10$ data sets). A score that disagrees with
the two tail scores flags an effect that lives in the bulk, not the tail.}
\label{tab:bulk-vs-tail}
\small
\begin{tabular}{@{}ll ccc@{}}
\toprule
Arm & Contrast & Brier (tail) & twCRPS (tail) & plain CRPS \\
\midrule
\multicolumn{5}{@{}l}{\emph{Ozone --- interpolation}}\\
Ozone & AR(2)$-$AR(1) & ns {\scriptsize $p{=}0.753$} & ns {\scriptsize $p{=}0.981$} & \textbf{worse} {\scriptsize $p{=}0.026$} \\
\addlinespace
\multicolumn{5}{@{}l}{\emph{Simulation spatial hold-out --- interpolation}}\\
Sim K{=}5 strong & AR(2)$-$AR(1) & ns {\scriptsize $p{=}0.221$} & ns {\scriptsize $p{=}0.979$} & ns {\scriptsize $p{=}0.581$} \\
\addlinespace
\multicolumn{5}{@{}l}{\emph{Time-block near-unit-root, leads 1--5 --- extrapolation}}\\
Time-block & AR(2)$-$iid & ns {\scriptsize $p{=}0.541$} & ns {\scriptsize $p{=}0.459$} & \textbf{better} {\scriptsize $p{=}0.026$} \\
Time-block & AR(2)$-$AR(1) & ns {\scriptsize $p{=}0.695$} & ns {\scriptsize $p{=}0.459$} & \textbf{better} {\scriptsize $p{=}0.033$} \\
\bottomrule
\end{tabular}
\end{table}


\section*{裁決：所有 AR(2) 效應都住在 bulk，尾部一律測不到}

三個 arm 的模式\textbf{完全一致}：

\begin{itemize}
  \item \textbf{臭氧（內插）}：plain CRPS 顯示 AR(2) 較差，但兩個尾部分數皆 ns。
  \item \textbf{模擬（內插）}：三個分數全 ns——徹底的虛無，無論 bulk 或 tail。
  \item \textbf{time-block（外推）}：\textbf{這是最關鍵、也最意外的一格}。plain
        CRPS 在近單位根、leads 1--5 顯示 AR(2) 顯著\emph{較佳}（先前 folder 04/05
        的亮點）；但一旦約束到尾部帶，twCRPS 的 $\Delta$ 掉回 $\approx 0$
        （$p_W\approx0.46$），單門檻 Brier 也 ns。
\end{itemize}

換言之，\textbf{plain CRPS 偵測到的所有 AR(2) 效應（臭氧的「較差」、time-block 的
「較佳」）都是分布主體（bulk）的現象，沒有一個發生在門檻超標的尾部。}

\paragraph{物理解讀（time-block）。}AR(2) 靠持續性改善的是\emph{未來日的水準／
中央預測}——plain CRPS 大力獎勵這一點；但「某站某日是否突破高門檻」由 AR(2)
無法更準預測的創新噪音主導。所以增益落在 level，不在 extremes。

\section*{對論文的（誠實）修正}

這修正了 folder 04/05 的措辭。先前把 time-block 的 plain-CRPS 增益當作
「AR(2) 是外推的預測工具」的實證支撐；尾部檢定顯示，\textbf{對本論文真正的目標
（門檻超標預測），連 time-block 的 AR(2) 增益也不成立}——它是 bulk 增益。

因此最終、最誠實的結論是：

\begin{quote}
用論文真正在乎的\textbf{尾部}分數衡量（單門檻 Brier，以及檢定力更高的尾部帶
twCRPS），\textbf{AR(2) 在所有情境——內插與外推——都不改變門檻超標的預測技巧}。
plain CRPS 觀察到的差異是分布主體的效應，與超標目標無關。
\end{quote}

這比原本「AR(2) 是外推工具」更保守，但證據力更強：null 現在由\textbf{兩個}尾部
分數、跨\textbf{三個} arm 一致支撐。它也讓論文的定位更乾淨——兩項延伸都不改善
尾部預測，其價值在於把 STP 擴為可由\emph{尾部}分數資料驅動選擇的候選類，而非
提供穩定增益。

\medskip
\noindent\footnotesize 程式：\verb|ozone/.../twcrps_bootstrap.R|、
\verb|simstudy/ar2_twcrps_export.R|、
\verb|block1_positive_control/timeblock_twcrps_export.R|。所有 twCRPS 皆為
尾部帶 $[q_{.90},q_{.99}]$、$z$-空間 $G{=}31$ 網格梯形積分。
\end{document}
